\chapter{Gyakorlati alkalmazások: wiki és CMS}

A következő két minta ugyanazokra az alapokra épül, de eltérő termékigényeket mutat. A minimális wiki a stabil slugot, Markdown tartalmat és autentikált szerkesztést kapcsolja össze; az érettebb wikihez később optimistic locking és audit trail adható. A CMS-minta ehhez többfájlos projektstruktúrát, szerkesztői jogosultságot és public cache-t ad, majd külön felsoroljuk a production CMS további rétegeit.

A wiki minimális, fordítási fixture-ként is ellenőrzött változata az \code{examples/wiki/main.rw}; a CMS többfájlos mintája az \code{examples/news-site/}. Így az alábbi két alkalmazásminta nem csak könyvbeli pszeudokód.

\section{Egyszerű wiki}

Egy wiki minimális domainje oldalból áll: stabil azonosító, slug, cím és Markdown törzs. Az alábbi minta a nyelvi elemek összekapcsolását mutatja.

\subsection*{Modell}

\begin{lstlisting}[language=RWLang]
model WikiPage {
    id: Int
    slug: Slug
    title: String
    body: String
}
\end{lstlisting}

\subsection*{Lekérdezés}

\begin{lstlisting}[language=RWLang]
query fn wikiBySlug(db: Db, slug: Slug)
    -> Result<WikiPage, DbError> sql {
    SELECT id, slug, title, body
    FROM wiki_pages
    WHERE slug = :slug
}
\end{lstlisting}

\subsection*{Publikus oldal}

\begin{lstlisting}[language=RWLang]
page fn wikiShow(ctx: PageContext, db: Db, slug: Slug)
    -> Result<Html, PageError> {
    let page = wikiBySlug(db, slug)?;
    return Ok(html {
        @layout(Main, page.title) {
            <article>
                <h1>{{ page.title }}</h1>
                @markdown(page.body)
            </article>
        }
    });
}

route wikiShow GET "/wiki/:slug<Slug>" => wikiShow;
\end{lstlisting}

\subsection*{Szerkesztő form}

\begin{lstlisting}[language=RWLang]
form WikiForm {
    title<String>
    body<String>
    validate title length 2 200 body length 1 300000
}
\end{lstlisting}

A módosító route legyen autentikált, például \code{auth role Editor}. A write query \code{Transaction} capabilityt kapjon. A cím alapján generált slugot ne írd át vakon minden szerkesztéskor, ha a régi URL-ek stabilitása fontos.

\subsection*{Canonical slug és alias}

Érettebb wikiben külön slug-alias/history tábla tartja a régi címeket. A runtime canonical route-ból tud 301-es redirectet készíteni, így nem kell tetszőleges redirect URL-t építeni.

\subsection*{Mit adj hozzá egy valódi wikihez?}

\begin{itemize}
  \item optimistic locking a konkurens szerkesztésekhez;
  \item üzleti audit trail a módosításokhoz;
  \item objektumszintű authorization, ha nem minden oldal azonos jogosultságú;
  \item backup/restore policy az adatbázisra és az esetleges AppFs adatra.
\end{itemize}


\section{CMS és híroldal}

A repository \code{examples/news-site/} példája már egy kis CMS magját mutatja: model, typed SQL, layout, Markdown, szerkesztői action és public cache.

\subsection*{Projektstruktúra}

\begin{lstlisting}
main.rw
models.rw
queries.rw
components.rw
pages.rw
actions.rw
\end{lstlisting}

\subsection*{Cikkmodell}

\begin{lstlisting}[language=RWLang]
model Article {
    id: Int
    slug: Slug
    title: String
    body: String
}
\end{lstlisting}

A szerkesztői input külön named form-sémát kap:

\begin{lstlisting}[language=RWLang]
form ArticleForm {
    title<String>
    body<String>
    validate title length 2 200 body length 1 300000
}
\end{lstlisting}

\subsection*{Publikus cikkoldal}

\begin{lstlisting}[language=RWLang]
page fn article(ctx: PageContext, db: Db, slug: Slug)
    -> Result<Html, PageError> {
    let article = articleBySlug(db, slug)?;
    return Ok(html {
        @layout(Main, article.title) {
            <article>
                <h1>{{ article.title }}</h1>
                @markdown(article.body)
            </article>
        }
    });
}

route article GET "/cikk/:slug<Slug>"
    cache public ttl 60
    => article;
\end{lstlisting}

Public cache csak olyan oldalra való, amely nem session- vagy userfüggő.

\subsection*{Szerkesztői létrehozás}

\begin{lstlisting}[language=RWLang]
action fn create(ctx: ActionContext, db: Db,
                 title: String, body: String)
    -> Result<Redirect, PageError> {
    let articleSlug = slug(title);
    transaction db {
        insertArticle(tx, articleSlug, title, body)?;
    }
    return Ok(redirect("/admin/articles"));
}

route articleCreate POST "/admin/articles"
    form ArticleForm
    auth role Editor
    invalidate cache article
    => create;
\end{lstlisting}

A fenti invalidáció route-szintű generationváltás: nem csak az új slug egyetlen cache-kulcsát törli, hanem az \code{article} route korábbi generationjét teszi elérhetetlenné. Egyszerű mintában ez helyes és kiszámítható, de nagy forgalmú CMS-nél érdemes tudatosan mérlegelni, mely publikus route-okat kell valóban invalidálni egy adott módosítás után.

\subsection*{CMS-ből production rendszer}

A következő lépések adják a gyakorlati production CMS gerincét:

\begin{itemize}
  \item publikálási és domain-invariánsok külön validációja;
  \item draft/published állapot enumként;
  \item canonical slug history;
  \item optimistic locking szerkesztéskor;
  \item business audit trail;
  \item media library és AppFs a képekhez;
  \item role-based szerkesztő/publisher policy;
  \item rate limit és resource profile a drágább route-okra;
  \item backup, restore drill és kontrollált deployment.
\end{itemize}

A CMS-ben az eredeti Markdown legyen a source of truth; ne tárolj alkalmazásból származó raw renderelt HTML-t.
